\section{Manipulación de arrays y filtrado}

Además del \textit{.reshape} visto en la sección de propiedades y la indexación y \textit{slicing}, existen otras formas de manipular y filtrar datos para que confiesen lo que necesitamos.

\subsection{Manipulación de arrays}

\subsubsection{Concatenación y apilamiento}

Si usted, como todo ser normal, viene del mundo de las hojas de cálculo, la concatenación de \textit{arrays} es simplemente la forma técnica de decir "copiar y pegar rangos de celdas" para armar una tabla más grande.

\paragraph{Concatenación (\textit{np.concatenate})}

Es la forma más general de unir dos o más arreglos. Esta función pega una lista de \textit{arrays} uno tras otro a lo largo de un eje específico.

Sin embargo, el autor sabe lo difícil que puede ser visualizar ejes en la cabeza sin que le empiece a doler la nuca, así que lo vamos a explicar en ``modo Umpa Lumpa''.

\begin{ejemplo}
Imagine que se encuentra en la fábrica de chocolates de Willy Wonka (esperemos que Paramount no me cobre derechos de autor). Usted y su pequeño compañero de diminuto tamaño están fabricando dos tipos de dulces: los rojos y los azules (dos arreglos vectoriales).
$$
\begin{aligned}
\text{\textit{Caramelos rojos}} &= [R, R, R] \
\text{\textit{Caramelos azules}} &= [A, A]
\end{aligned}
$$
Wonka grita: ¡Concatenen! Lo que usted hace es colocar los caramelos azules justo donde terminan los rojos, dándonos una sola fila:
\[
\left[R, R, R, A, A\right]
\]
\end{ejemplo}

Pero, mi pequeño compañero Umpa Lumpa, a veces la cosa se complica cuando dejamos los caramelos y pasamos a las bandejas (matrices) de \textit{brownies}.

Si utilizamos \textit{np.concatenate()} en matrices, debemos definir los \textit{axis} (ejes). El \textit{axis=0} hace que coloque una bandeja encima de la otra. Básicamente, la pila crece hacia arriba.

Pero tenga cuidado: apilar \textit{brownies} no es una actividad fácil cuando hablamos de millones de bandejas. Todas deben tener exactamente la misma dimensión en el ancho. No podemos poner una bandeja de cuatro \textit{brownies} sobre una de dos, porque se nos caería la torre y Wonka nos despediría sin indemnización.

\begin{nota}
    El \textit{axis=0} afecta a las filas. Si concatenamos dos matrices de \(2 \times 3\) usando \textit{axis=0}, obtendremos una de \(4\times 3\). El ancho se mantiene igual, pero la estructura ahora es más alta.
\end{nota}

``¿Qué pasa si en vez de ordenar por filas ordenamos por columnas?''. Pues lo único que debe hacer es utilizar \textit{axis=1}. Con esto le decimos a la máquina que pegue una bandeja al lado de la otra. Aquí también hay una condición: a Wonka (\textit{NumPy}) no le gusta que las bandejas tengan distinto largo (filas). Si intenta pegar una bandeja de 3 filas al lado de una de 2, la línea no encajaría y terminaría en la calle.

Véalo en el código de la fábrica:

\begin{pycode}{Concatenación}
import numpy as np

# Inventario de dulces
bandeja_1 = np.array([[1, 1], [1, 1]]) # 2x2
bandeja_2 = np.array([[2, 2], [2, 2]]) # 2x2

# Unir verticalmente (Hacer la torre más alta)
torre = np.concatenate((bandeja_1, bandeja_2), axis=0) # Resultado: Una matriz de 4x2

# Unir horizontalmente (Hacer la fila más larga)
fila_larga = np.concatenate((bandeja_1, bandeja_2), axis=1) # Resultado: Una matriz de 2x4
\end{pycode}

\begin{nota}
    \textit{Truco de sintaxis:} Fíjese que los arreglos que quiere unir van dentro de un par de paréntesis extra \texttt{((a, b))}. Esto es porque \textit{np.concatenate} espera una "secuencia" (como una tupla o lista) de \textit{arrays}, no los \textit{arrays} sueltos.
\end{nota}

Ahora usted entiende que es cocatenar una matriz y un vector (sin la parte matemática). 

Pero ¿Para qué sirve esto fuera de la fábrica de chocolates? Usted, que probablemente no tiene un contrato vigente con el señor Wonka, se preguntará: ``¿Y para qué quiero saber esto?''. La respuesta es que, en el análisis de datos, la concatenación es su mejor amiga cuando los registros vienen por separado.

Imagine estas situaciones cotidianas para un analista: 
\begin{itemize}
    \item Suponga que su sistema genera un registro de ventas diferente cada día. Si usted quiere realizar un análisis semanal, no va a mirar siete archivos por separado como un monje medieval; usted va a concatenar esos siete registros para tener un reporte semanal sólido en un solo \textit{array}.
    \item \textit{Fusión de fuentes:} Si tiene dos bases de datos que hablan del mismo tema (por ejemplo, encuestas de dos sucursales distintas), al concatenarlas fortalece su muestra. Una muestra más grande significa una estadística más precisa y una media de datos que no miente tanto.
\end{itemize}

En resumen, concatenamos para unificar. Un dato suelto es una anécdota, mil datos son una estadística.


\paragraph{Apilamiento}

Seamos sinceros, estar peleandose con los ejes (\textit{axis=0, axis=1}) es una actividad de alto riesgo para su salud mental. \textit{NumPy} sabe que usted es un ser humano (o algo parecido) y le ofrece ``atajos'' semánticos que hacen lo mnismo que concatenate, pero con nombres que hasta un niño de cinco años entenderia.

Si \textit{concatenate} es la herramienta multiuso, \textit{vstack} y \textit{hstack} son herramientas especializadas. Sus nombres vienen del inglés \textit{vertical stack} (apilamiento vertical) y \textit{horizontal stack} (apilamiento horizontal).

\begin{itemize}
\item \textit{np.vstack((a, b)):} Es el equivalente a \textit{concatenate} con \textit{axis=0}. Apila los arreglos verticalmente, uno sobre otro, como si estuviera armando una torre de panqueques.
\item \textit{np.hstack((a, b)):} Es el equivalente a \textit{concatenate} con \textit{axis=1}. Apila los arreglos horizontalmente, uno al lado del otro, como si estuviera formando una fila en el supermercado.
\end{itemize}

Fíjese qué limpio queda el código cuando dejamos de hablar de "ejes" y empezamos a hablar de "direcciones":

\begin{pycode}{Atajos de apilamiento}
import numpy as np

a = np.array([1, 1, 1])
b = np.array([2, 2, 2])

# Apilamos hacia arriba (Vertical)
vertical = np.vstack((a, b)) # Resultado: una matriz de 2x3

# Apilamos hacia el lado (Horizontal)
horizontal = np.hstack((a, b)) # Resultado: un vector largo de 6 elementos

print("Vertical:\n", vertical)
print("Horizontal:", horizontal)
\end{pycode}

\begin{nota}
    \textit{Un detalle de vital importancia:} Al igual que con \textit{concatenate}, estas funciones esperan recibir una tupla (es decir, doble paréntesis \texttt{((a, b))}). Si olvida el paréntesis interno, NumPy pensará que le está pasando argumentos extra y le lanzará un error que le arruinará la tarde.
\end{nota}




\subsubsection{División de arrays (\textit{split})}

Si concatenar y apilar era el arte de construir bloques, \textit{np.split()} es la herramienta para desarmarlos. Esta función corta un \textit{array} en partes exactamente iguales y genera una lista de sub-arreglos más pequeños.

A diferencia del \textit{slicing}, que es como sacar una rebanada de pan, \textit{split} está diseñado para repartir todo el bloque de un solo comando.

\begin{nota}
    \textit{La regla de la división exacta:} \textit{np.split} es extremadamente estricto. Si usted le pide dividir un \textit{array} en tres partes, el tamaño total del arreglo debe ser divisible por 3. Si la división no da un número entero, \textit{NumPy} entrará en pánico y lanzará un error porque no sabe qué hacer con el elemento que sobra. 
\end{nota}

\textit{np.split()} tiene el siguiente modo de operar:
\begin{verbatim}
    np.split(array, cantidad de divisones)
\end{verbatim}

Veamos cómo se descuartiza un vector en el código:

\begin{pycode}{Uso de np.split}
import numpy as np

# Creamos un vector de 6 elementos
datos = np.array([10, 20, 30, 40, 50, 60])

# Dividir en 3 partes iguales
pedazos = np.split(datos, 3)

print(pedazos[0]) # [10, 20]
print(pedazos[1]) # [30, 40]
print(pedazos[2]) # [50, 60]
\end{pycode}

En el caso de las matrices, volvemos a encontrarnos con nuestros viejos amigos, los \textit{axis} (ejes), para elegir la dirección del "hachazo":
\begin{itemize}
\item \textit{axis=0}: Es un corte horizontal. Divide la matriz a través de sus filas (obtenemos trozos de arriba y abajo).
\item \textit{axis=1}: Es un corte vertical. Divide la matriz a través de sus columnas (obtenemos trozos de izquierda y derecha).
\end{itemize}

Vea el siguiente ejemplo:
\begin{pycode}{Split en Matrices}
import numpy as np

matriz = np.arange(16).reshape(4, 4)

# Dividir la matriz a la mitad (horizontalmente)
arriba, abajo = np.split(matriz, 2, axis=0) # Nos quedan dos matrices de 2x4

# Dividir la matriz a la mitad (verticalmente)
izquierda, derecha = np.split(matriz, 2, axis=1) # Nos quedan dos matrices de 4x2

print(arriba, abajo)
print()
print(izquierda, derecha)

\end{pycode}

¿Para qué sirve esto en la vida real? Usted usará \textit{np.split()} principalmente para dos cosas:
\begin{enumerate}
\item \textit{Separación de variables:} Si tiene una tabla donde las primeras columnas son características (como altura y peso) y la última es el precio, puede usar \textit{split} para separarlas y procesarlas por rutas distintas.
\item \textit{Procesamiento por lotes (\textit{batches}):} Si tiene un millón de registros y su computadora empieza a echar humo, puede dividirlos en 10 partes iguales y procesar una por vez para darle un respiro al procesador.
\end{enumerate}

\begin{nota}
Para facilitarnos la vida una vez más, \textit{NumPy} también incluyó las funciones \textit{np.vsplit} (vertical) y \textit{np.hsplit} (horizontal), que funcionan exactamente igual que los atajos de apilamiento, evitándole tener que escribir el \textit{axis} manualmente.
\end{nota}



\subsubsection{Transposición y permutación de ejes}

\paragraph{Transposición}

Transponer una matriz consiste en el arte de intercambiar sus filas por sus columnas. Si usted tenía una fila, ahora tiene una columna; si tenía una columna, ahora es una fila. Es, básicamente, rotar su tabla de datos 90 grados y darle la vuelta.

Matemáticamente, si un elemento estaba en la posición \(i,j\) donde \(i\) es una fila y \(j\) es una columna, la matriz transpuesta ahora sera \(j,i\), Esto se visualiza de la siguiente manera:
\[
A = \begin{bmatrix}
    1 & 2 \\
    3 & 4 \\
    5 & 6
\end{bmatrix}
\quad
\text{Transpuesta}
\quad
A^T= \begin{bmatrix}
    1 & 3 & 5 \\
    2 & 4 & 6
\end{bmatrix}
\]

¿De qué sirve transponer una matriz en la vida real?
\begin{itemize}
\item \textit{Reorientar datos:} Si sus sensores registraron datos en columnas pero su algoritmo los espera en filas, la transposición lo arregla en un milisegundo.
\item \textit{Álgebra Lineal:} Es fundamental para cumplir las condiciones de la multiplicación de matrices (que veremos si el destino lo permite).
\item \textit{Correlaciones:} Permite comparar cómo se relacionan diferentes variables entre sí al alinear correctamente los ejes.
\end{itemize}

Para transponer, \textit{NumPy} ofrece la función \texttt{np.transpose(array)}, pero como usted es un programador que valora su tiempo, usará la propiedad \texttt{.T}. Es corta, es elegante y le ahorra preciosos segundos de vida que puede gastar en ver videos de gatitos o tomar café.

Véalo en el código:
\begin{pycode}{Transposicion con .T}
import numpy as np

# Creamos una matriz de 2 filas y 3 columnas
matriz = np.array([
[1, 2, 3],
[4, 5, 6]
])

print(f"Forma original: {matriz.shape}") # (2, 3)

# Transponemos la matriz
matriz_t = matriz.T

print(matriz_t)
"""
Resultado:
[[1, 4],
[2, 5],
[3, 6]]
"""
print(f"Nueva forma: {matriz_t.shape}") # (3, 2)
\end{pycode}

\begin{nota}
    Si intentas transponer un vector de una sola dimensión (un array simple como \texttt{[1, 2, 3]}), verás que \textit{no pasa nada}. Esto es porque NumPy considera que un vector 1D no tiene filas ni columnas que intercambiar. Para "parar" un vector, primero deberías convertirlo en una matriz de \(1 \times n\) usando \texttt{.reshape()}.
\end{nota}



\paragraph{Permutación de ejes}

Si la transposición (\texttt{.T}) es como dar vuelta toda la media, \texttt{np.swapaxes()} es un movimiento mucho más preciso: intercambia exactamente dos ejes entre sí, dejando todos los demás intactos en su posición original. Su estructura consiste en:
\begin{verbatim}
    np.swapaxes(array, eje_1, eje_2)
\end{verbatim}


Esta función es la favorita de quienes trabajan con datos multidimensionales, como imágenes o tensores complejos, donde a veces solo necesita rotar dos dimensiones específicas (por ejemplo, intercambiar el ancho por el alto) sin alterar la profundidad o el número de canales de color.

En una matriz simple de dos dimensiones, usar \textit{swapaxes} es exactamente lo mismo que usar \textit{.T}. La verdadera diferencia —y donde usted agradecerá su existencia— aparece cuando trabajamos con tres o más dimensiones.

Vea este ejemplo:

\begin{pycode}{transpose vs swapaxes}
import numpy as np

# Creamos un "cubo" de datos (2 matrices de 3x4)
cubo = np.zeros((2, 3, 4)) # np.shape: (0: profundidad, 1: filas, 2: columnas)

# 1. Usando transpose para invertir TODO
cubo_t = np.transpose(cubo) # La forma pasará a ser (4, 3, 2)

# 2. Usando swapaxes para cambiar solo FILAS por COLUMNAS
cubo_swapped = np.swapaxes(cubo, 1, 2) # La forma pasará a ser (2, 4, 3)
# Solo tocamos el eje 1 y el eje 2. La profundidad (eje 0) queda igual.

print(f"Original: {cubo.shape}")
print(f"Transpose: {cubo_t.shape}")
print(f"Swapaxes: {cubo_swapped.shape}")
\end{pycode}

Como puede observar, \texttt{swapaxes} es mucho más respetuoso con la estructura que usted no quiere tocar. Mientras que \texttt{transpose} lo pone todo patas arriba.




\subsection{filtrado de datos}

Además de la indexación y la indexación booleana existe el filtrado condicional con \texttt{np.where()}. Esta función es, en esencia, un bloque \textit{if-else} convertido en un superatleta vectorizado.

Su estructura consiste en:
\begin{verbatim}
    np.where(condición, valor_verdadero, valor_falso).
\end{verbatim}

Suponga que quiere evaluar una lista de alumnos para separar el trigo de la paja (aprobados de desaprobados). En lugar de ir uno por uno preguntando la nota, usted le da la orden a NumPy y él se encarga del resto:
\begin{pycode}{Uso basico de np.where}
import numpy as np

notas = np.array([3, 8, 5, 2, 10])

# Si la nota es >= 4, ponemos 1 (Aprobado), sino 0 (Desaprobado)
resultados = np.where(notas >= 4, 1, 0)

print(resultados) #Resultado: [0, 1, 1, 0, 1]
\end{pycode}

Existen dos razones fundamentales por las que usted, si aprecia su tiempo, usará \texttt{np.where} en lugar de un \texttt{if} tradicional de \textit{Python}:
\begin{itemize}
\item \textit{Vectorización (Velocidad):} Como mencionamos antes, NumPy procesa todos los elementos al mismo tiempo. Si tuviera un millón de filas, \textit{np.where} tardaría milisegundos, mientras que un bucle \textit{for} con un \textit{if} le daría tiempo de ir a prepararse un café (y quizás hasta de cultivarlo).
\item \textit{Limpieza de Datos:} Es la herramienta ideal para corregir errores absurdos en sus \textit{datasets}. Por ejemplo: "Si el precio es negativo, póngale 0; de lo contrario, deje el precio original".
\end{itemize}

Vea cómo limpiar su honor (y sus datos) con un solo comando:
\begin{pycode}{Limpieza de datos con np.where}

precios = np.array([100, -50, 200, -10])

# Si el precio es menor a 0, lo corrijo a 0. Si no, mantengo el precio original.
precios_limpios = np.where(precios < 0, 0, precios)

print(precios_limpios) # Resultado: [100, 0, 200, 0]
\end{pycode}

Si usted le pasa a \texttt{np.where} solo la condición (sin los valores de verdadero/falso), la función cambia de personalidad: en lugar de devolver valores transformados, le devolverá los índices exactos donde se cumple lo que usted pidió.

\begin{pycode}{np.where para encontrar indices}
import numpy as np

edades = np.array([10, 25, 18, 40])

indices = np.where(edades >= 18) # Resultado: (array([1, 2, 3]),)
print(indices)
\end{pycode}

\begin{nota}
Fíjese que el resultado es una tupla que contiene un \textit{array}. Esto parece innecesariamente complejo, pero es así porque está preparado para matrices de muchas dimensiones donde usted recibiría una lista de índices por cada eje. Para nosotros, por ahora, basta con saber que ahí están nuestras coordenadas.
\end{nota}